\section{Open Questions}

This replication exercised the paper's central analytical claim (permutation invariance of GM-QAOA), its Eq.~(8) closed form, and its X-vs-GM introductory advantage claim on a single structured MAX-CUT instance. Five concrete follow-ups remain open:

\begin{enumerate}
  \item \textbf{Numeric fit of the main asymptotic theorem (C5).} Does the paper's exponential-in-$n$ round-count lower bound for GM-QAOA on complete-bipartite MAX-CUT hold quantitatively when tested by a statevector scaling sweep up to $n \approx 14$, and does the fitted log-slope match the theorem's constant?
  \emph{Basis:} C5 is the main new theorem but was not exercised here --- a single small-instance run cannot verify an asymptotic bound.
  \emph{Next steps:} implement GM-QAOA on $K_{a,b}$ for $(a,b) \in \{(2,2),\ldots,(7,7)\}$, at each $n$ find minimum depth $p^\ast(n)$ needed to reach approximation ratio $0.9$ (COBYLA, 100 restarts), and plot $\log p^\ast(n)$ vs.\ $n$.

  \item \textbf{Graph-family dependence of the X-vs-GM gap.} Does the $+0.068$ to $+0.147$ gap observed on one 6-node graph generalize across graph families, or is it graph-topology-dependent?
  \emph{Basis:} C4 was exercised on a single graph. The paper cites this as a general structured-problem phenomenon but one instance is not a distribution.
  \emph{Next steps:} sweep 30 random 3-regular graphs and 30 $G(n{=}8,p{=}0.5)$ graphs at depths $p=1,2,3$; report the distribution of (X-ratio minus GM-ratio) across families.

  \item \textbf{Warm-starting GM-QAOA vs.\ the X-mixer.} Does warm-starting GM-QAOA (initialize state to a classical relaxation solution rather than $|s\rangle$) recover any of the approximation-ratio gap, or does structural blindness still dominate?
  \emph{Basis:} the Grover mixer's reflection about $|s\rangle$ is intrinsic to the permutation-invariance identity. Warm-starting breaks that assumption, so it is a natural probe of the paper's structural-blindness argument.
  \emph{Next steps:} replace $|s\rangle$ with a state biased toward a Goemans--Williamson SDP rounded cut; rerun the C4 comparison on Graph A and on 20 random 3-regular graphs.

  \item \textbf{Behavior on constrained problems (MaxIndependentSet, MaxCoverage).} How does GM-QAOA compare to the X-mixer on constrained problems where the Grover mixer's ability to stay in a feasible subspace is often cited as its main advantage?
  \emph{Basis:} the paper's negative result is on unconstrained MAX-CUT. Grover-style mixers were originally motivated (Hadfield et al.~2019) precisely for constrained problems, a regime not touched here.
  \emph{Next steps:} implement MaxIndependentSet on $n=6,7,8$ with a Grover mixer projecting to the feasible subspace; compare against X-mixer QAOA with a soft infeasibility penalty; report approximation ratio and feasibility fraction.

  \item \textbf{Noise robustness of the permutation-invariance identity.} Under a realistic depolarizing / amplitude-damping noise model at $10^{-3}$ per two-qubit gate, does C1 still hold approximately, or does noise break the symmetry and change the qualitative X-vs-GM comparison?
  \emph{Basis:} every result here is noiseless statevector. The paper's structural-blindness argument is a strict identity of the unitary operator; noisy channels are not permutation-symmetric in general.
  \emph{Next steps:} use Qiskit Aer density-matrix simulator with per-gate depolarizing $p \in \{10^{-3},10^{-2}\}$; rerun C1 on Graph A at $p_{\mathrm{depth}}=2$ and rerun C4 under noise; report whether X-mixer still beats GM-QAOA.
\end{enumerate}
